消元已经能回答一个方阵是不是可逆，可它交出来的是一句是非判断。有些问题要的是量：这个线性变换把一块区域撑大了还是压小了，倍数是多少；它有没有把空间翻过一面；如果某个方向被彻底压平，这件事会以什么形式暴露出来。

本章把这三个问题的答案塞进同一个标量。\(\det A\) 的绝对值是体积的缩放因子，符号记录定向有没有翻转，取值为零说明至少一个维度塌了。压缩得这么彻底，好处是判断可逆只需看一个数；代价是方向信息全部丢失——被压扁的是哪个方向、离奇异还有多远，行列式一概不答。

\begin{center}
\begin{tikzpicture}[x=1cm,y=1cm,font=\footnotesize,>=stealth]
  \draw[black!55,fill=black!4] (0,-0.5) rectangle (1.8,0.5);
  \node at (0.9,0) {方阵 \(A\)};
  \draw[->,black!60] (1.9,0) -- (2.7,0);
  \draw[black!55,fill=blue!12] (2.8,-0.5) rectangle (4.4,0.5);
  \node at (3.6,0) {\(\det A\)};
  \draw[->,black!60] (4.5,0.25) -- (5.4,1.4);
  \draw[->,black!60] (4.5,0) -- (5.4,0);
  \draw[->,black!60] (4.5,-0.25) -- (5.4,-1.4);
  \draw[black!45] (5.5,1.1) rectangle (9.9,1.9);
  \node at (7.7,1.5) {绝对值：体积缩放了几倍};
  \draw[black!45] (5.5,-0.4) rectangle (9.9,0.4);
  \node at (7.7,0) {符号：定向有没有翻过来};
  \draw[black!45] (5.5,-1.9) rectangle (9.9,-1.1);
  \node at (7.7,-1.5) {取零：至少一个维度塌了};
  \draw[->,dashed,black!50] (3.6,-0.6) -- (3.6,-2.3);
  \node[align=center,black!65] at (3.6,-2.75) {被压扁的是哪个方向、离奇异多远\\这些留给特征值与奇异值};
\end{tikzpicture}

\emph{一个标量能承担的和不能承担的。左侧三条是行列式给出的判断，底下那条虚线提醒它把什么丢掉了。}
\end{center}

支撑全部内容的只有三条要求：对每一列线性、两列相同则取零、单位矩阵取一。\hyperref[sec:03-Linear-Algebra/04-Determinants/01]{``行列式与体积：线性变换是放大、翻转，还是压扁空间''}说明这三条为什么就是``有向体积''的全部内容，以及交换变号、相关列取零、行列式为零等价于奇异这些常被当作条目背诵的规则，怎样从它们直接掉出来；这一篇也澄清一处高频误读——行列式很小并不代表矩阵病态。\hyperref[sec:03-Linear-Algebra/04-Determinants/02]{``行列式的计算规则与推论''}转向计算与后果：消元怎样一路记账把行列式算出来，\(\det(AB)=\det A\det B\) 为什么一句几何话就说完，余子式展开和 Cramer 法则为什么在理论上好用而在数值上无人问津，以及这个体积因子怎样以 Jacobian 行列式和 \(\log\det\) 的形式出现在概率密度变换、normalizing flow 和高斯似然里。

两篇按顺序读，第一篇是第二篇的全部依据。动手之前需要熟悉消元与矩阵乘法；若对``某个方向被压扁''这句话还没有画面，可以先回看 \hyperref[sec:03-Linear-Algebra/02-Vector-Spaces/03]{``矩阵的四个基本子空间：输入、输出与丢失的信息''}。读完之后，看到一个矩阵你应当能说出它对体积做了什么、对定向做了什么，也能说出在什么场合不该拿行列式当作数值健康度的指标。
